\section{Discussion and Limitations}

The sentinel works because configuration semantics identify a high-value measurement target. A global historical residual mixes heterogeneous delay regimes, whereas the candidate OSPF change identifies the direct link whose state determines the post-change path. Measuring the operational/twin discrepancy on that link removes most of the dominant prediction error before deployment and remains available even for candidates that are rejected.

The result suggests a broader control pattern: identify resources activated by a proposed change, obtain targeted operational measurements, compare them with the twin representation, correct or invalidate the twin prediction, and abstain when the corrected score violates policy. This complements, rather than replaces, reachability verification and consistent-update mechanisms.

External validity is limited. The topology has three routers, one OSPF area, one candidate link, and one weight change; drift is fixed latency on that link. We do not evaluate capacity reduction, loss, queue buildup, multi-link or topology drift, failed links, multi-area or multi-vendor networks, concurrent controllers, or production scale. Mean ICMP RTT is a narrow service metric. Calibration (22) and holdout (33) are small, Wilson intervals remain wide, and exchangeability does not hold well enough to preserve nominal conformal coverage. The harness also deploys rejected candidates offline to obtain labels, so the artifact evaluates a gate rather than a live enforcement controller. These limitations motivate larger topologies, multiple candidate links and impairments, explicit runtime measurements, and adaptive or risk-controlling calibration under nonstationarity.
